I have a little story to tell -- a childhood story and one which I tell to the children. When I was little, perhaps eight years old, I used to herd sheep up on top of Lukachukai Mountain with a cousin. My cousin was just a little bit older than I was -- he was perhaps nine years old at that time. That was about how big we were when we herded sheep by ourselves.

At that time up there on the mountain there were those that live in mountains, bears as they are called. But one does not mention their name ``bear''. People who live over there toward the mountains never say ``bear''. [...]\footnote{The English translation in \citet[51]{Young:1954jz} adds the following sentence: ``But at the time when we were small and used to herd our sheep up there we did not know things like this.'' This does not match any Navajo counterpart.} Only my cousin had a smattering of such information, and he used to teach me what he knew. ``There are certain things here in the mountains. There is something called `bear,' but one never refers to it by that name. People do not say `bear', so never say that word,'' He would say as he taught me. He was older than I was so he had learned some things. ``Bear is called by other names,'' he would say, instructing me. ``In the summertime it is called `bushy-tail'.'' Then in the morning when we would take out the sheep, he would again remind me of this lesson.

We would take the sheep out, and we would herd them in a real wilderness area. These areas were dense with pine, spruce, aspen and oak. Water flowed out from the canyons, and little lakes lay in the flats beyond. It was very beautiful up there on the mountain at that time. We had a large herd of sheep. When they spread out into the brush we would just let them go and follow them. My cousin would tell me to go a certain direction, and we would then separate. When the sheep would be on the side of the mountain, I'd skirt them from above. And when I would go up there, I'd remember what he had told me about not mentioning the word `bear'. I used to go around wondering why under the sun a person should not say the word `bear', so one day I decided to have some fun with my cousin by running down toward him as he walked along below me. So I dashed toward him. ``Here comes a bear, cousin. Look out, look out! A bear! A bear!'' I said, as I went zooming by him. After I had passed beyond him I turned around and came back to him.

I came back to where he was, and he was certainly angry. ``What do you mean by dashing around and saying that? Don't you remember that I told you that you are not say that under any circumstance? There could be some of them right here nearby. Don't you remember? That's why I told you that. I told you that they have a different name,'' he said as he scolded me.

Then we started to follow the sheep again. We trailed them for quite a while when suddenly, from within a nearby oak clump, there came the sound of dogs barking loudly. What is it that the dogs can be heard barking at? I wondered what the dogs were barking at so I ran over there. My cousin followed me. I got over to where the dog was barking, and there I saw a bear standing. I stopped right there and looked at him. He was really angry. He stood growling. He was picking up things that were scattered on the ground about him, such as sticks and stones, and threw them at the dogs. He used sticks to strike among them. The dogs ran around him barking fiercely. I got that glimpse of him and then I dashed away. My cousin too lost no time in running away. I stopped at a place far removed from the scene. It was only after I stopped that my cousin caught up with me. Then we started off again, and went over to where the sheep were walking in order to get in ahead of them and cut them off. We headed them off toward the flat. We herded them back down to the open country. We stayed down there and herded the sheep until noon. That is how I saw my first bear. I saw another one on the following day, but at a distance, and walking on the mountainside. It turns out that must be a bear.

In the fall of that same year we moved from the mountain back down into the valley to herd down there. We had the rams in the herd, and they used to really fight with one another. When the rams would attack and butt each other we would each get astride a ram and then urge them to attack. We would sit on them, and when their heads came crashing together we would really die laughing. Some times we would fall off, and we would really laugh then too. That is how we used to have fun with the rams. Sometimes the rams would get far apart and my cousin would step in between them and then jump aside as their heads came together. The sheep would be grazing nearby, and we would spend the whole day playing with the rams.

Once we had billy goats in the herd. That took place down on the plain too, not in the mountains. We had one billy goat in the herd that was just a little bit mean. If a person would bother him, like pinching him or walking behind him, he would get mad. We always herded sheep with a burro, sitting one behind the other. But when one would be riding the burro the billy goat would merely look up if one teased him. He wouldn't do anything under these circumstances even if one got very close to him. But the minute you dismounted in front of him, he came for you. That is how we used to herd sheep when this dangerous billy goat was present.

One time we got to wondering what he might do to us. He was grazing at a distance, and we went up to him on the burro. When we got there I dismounted right in front of him. I went up behind him, and he wheeled around to face me. Then he started trotting toward me, so I ran over behind the burro. Then my cousin, with whom I was herding sheep, jumped off the burro too. He approached the goat that was chasing me and it ran toward him. The goat went for him, but just before it could hit him he jumped outside and the goat went on past him. The goat went a little distance, turned around and came for us again. The faithful burro that we rode was now behind the billy goat. So we ran in the direction opposite the burro, and the goat trotted along behind us. He was driving us away from the burro. We kept right on running ahead of him, with him trotting along behind us. We were running in the direction of home, which lay about three miles away. We got to a little hill near home, on which there were some outcroppings of rock. There we crawled into a little niche in the rock. The goat came right up to the opening, as we lay there in the space beneath the rock. He took his spite out on the rock, butting it, and bit at our feet, He did every thing he could think of, in an effort to get us out, while we just lay there kicking at him. My cousin told me not to bother him, to just keep quiet and lie still, so we just lay still. The goat tired in his efforts to get us so he calmed down too, and just stood there watching us.

After a long time the goat went away. The sun had just gone down, and as soon as the goat went away we crawled back out of the hole. We could see him going along homeward. So we followed him back home. Fortunately all the sheep had returned, as well as the burro. We went about the task of milking goats for our supper. My cousin did the milking, while usually held the goat's head for him. That was my job. While we were doing our milking this old goat that had chased us all over the place got mad at us again. He suddenly came for us from a place close by. I saw him coming, but before I could turn he butted me in the hip and knocked me down. My legs went up in the air and one of my moccasins fell off and went flying. I leapt to my feet and started to run with only one moccasin on. My cousin nearly died laughing at what happened there.

When we were children we used to do all of these things. It would have been otherwise had our mothers been there to watch over us and tell us what we shouldn't do. We were just little boys alone, wandering around behind the sheep. We just did any mischievous thing we felt like doing. My mother and father lived far away on the farm.

These things that I have told are things that happened when I was six or seven years old. When I was perhaps eight years old the people in the whole area hoaxed one another. There was a tale that there was going to be a terrible flood in which the waters would carry everything before them, and a lot of the people believed it. So they moved up onto the mountain. From this day on it was about thirty-two years ago that it happened in the summer. It was said that the first word of this had come from up north. The rumor was carried everywhere. In one day it spread all over the country. ``The flood is going to be terrible. The ocean is overflowing. It will lie over us. The waters will come rolling over us. Go up to those high hills! Act quickly! Get up on those high places, for this is the greatest amount of water that ever lopped. There's no way of stopping the rise of this flood water. Hurry! Take all the livestock you can'', people said as this terrible story went around. I remember this event. We were living in the valley and I was still herding with my cousin when this story went around. At this time we were down at Many Farms with our sheep, and word was brought to my father and mother. They didn't believe it. ``This cannot be true. They're just talking and pulling our leg. How can the ocean go uphill? I don't believe that,'' said my father when the word was brought to him.

But all of our neighbors and relatives believed it, so there was really some movement. They started moving. They begged each other to go, saying, ``It's true. It's true. What's going to hold back the water? If it covers us over we'll certainly all die,'' the people said as they hastily got ready to move.

Everybody was moving in a different direction. Some moved over against Lukachukai Mountain. And some moved over to Black Mountain. This moving lasted for about three days. But the story had not affected us. Both my father and my mother had refused to believe it. We just kept on herding the sheep. People said the water was coming. When my cousin and I laid down to sleep at night we would wonder what the night might hold. We lived in a tent, for we stayed where the sheep were. But morning came and nothing had happened. So we spent another day herding and there was no water in sight. There were merely some black clouds in the distance. Perhaps there was rain here and there. Water wasn't expected throughout the following day. But no water came. That's what happened once. The people were victims of a flood hoax and all moved away. The old folks know about this. Those who moved away know about it. Every once in a while they make fun of one another about it. It's just a funny story. When they moved over to the mountains they went in wagons. They took their wagons just as close to the mountains as they possibly could. They made their wagons fast to pine trees and then continued on up the mountainside carrying their belongings on their backs. They fastened them to oak trees. To keep the water from carrying off their wagons people loaded them with rocks. They did this, thinking that they would still have their wagons when the flood receded. This event affected a large area but now people merely tell of it as a funny story.

One time the elder brother of this cousin I've been telling about and I were herding sheep alone and he told me, ``Cousin, I'm going to herd by myself today. You just stay here and keep an eye on the camp. You watch our things,'' he said. I said all right, so he set off with the sheep by himself. He made me stay at home to watch the place. My only companion was a dog, and we just stayed in a tent.

All of a sudden, about noon, some coyotes attacked me. At that time there was a great many coyotes. There were coyotes everywhere. From behind the tent a coyote sprang on me. The dog that I had there was a big white one with black ears. I was watching the place with a dog named Gogo when the coyote sprang on me from behind the tent. Just as he attacked me Gogo sprang to the rescue and drove him off me. I was frightened to death. I started running and the dog right after me. From the distance I kept my eye on our things. And then another coyote attacked me from another side. Just as this one jumped on me the dog came to my rescue, and drove that one off too. That is how the dog protected me. Now as I think about it, it is a good thing to have a dog. A dog can save one's life. I still remember my dog. I don't forget him, because he protected me. A dog will save a person when he is alone without protection. I found that out.

It was about a year after that a lot of people died. The white people say that it was in 1918. It was at that time that people died from an epidemic of influenza. In some places whole families died. In other places just a few survived. I had a hard time during that period. I used to have some brothers and sisters. There were nine of us. I was the ninth. Eight of us died. I just barely pulled through. I had a hard time during the flu epidemic. At that time there were no doctors nearby. Probably the only hospital then was the one at Fort Defiance, and possible the one at Ganado. Just very recently some hospitals were built, but we do not have enough of them. As a result, when people got sick they really suffered everywhere. Back then that was how I suffered among others. I have told the things that happened when I was a boy. I'm telling it to you children. Some of you children have mothers and fathers who are living. Some of your mothers are not living. But even so, don't forget your parents. I am saying so. Sometimes a child gets into mischief and doesn't mind his mother. Even so, he gets wiser and wiser. His mind grows and strengthens, and when he is told to do something he knows what to do. As time goes on, a child develops his thinking from his mother and father. That's the Navajo way. That's the way children grow up. His grandparents are like teachers to him. He asks them how The People used to live, and they tell him. It is akin to the textbooks of the white people. It is said that, in the past, the people really made it a point to teach one another. But today we are lax and do not teach our children much. Some of us neglect our children. We somewhat don't think for our children. It's well doable for us to teach our children those things. Even though your children are taught for you in school, we should not lose sight of our teaching. I am now telling these stories of mine to your children. Thes stories that you heard are those that I remember happening when I was a child. So then it's good that one recalls those stories, those things that happened. That's how you have memories. That's how you strengthen your mind.